\section{Discusión Crítica}
\label{sec:discusion}

\subsection{Contraste entre Teoría Asintótica y Desempeño Empírico}
La cota amortizada de Dijkstra con Fibonacci Heap mejora la dependencia respecto de $|E|$ frente a un heap binario:
\begin{equation}
    T_{\text{Fib}} = O(|E| + |V| \log |V|) \quad \text{frente a} \quad T_{\text{Bin}} = O((|V| + |E|) \log |V|)
    \label{eq:asintotica_comparativa}
\end{equation}
Esa diferencia asintótica importa cuando crece el número de aristas, pero no fija las constantes de una implementación. En las diez escalas sintéticas medidas en Python 3.12, el Fibonacci Heap registró medianas entre 1.05 y 1.59 veces las del binario.

El Fibonacci Heap manipula nodos enlazados y consolida árboles al extraer el mínimo. Es una posible fuente de sobrecosto, pero no se perfilaron asignaciones ni memoria; los tiempos y contadores de operaciones no permiten cuantificarla.

\subsection{Montículo 4-ario: cotas y mediciones}
El 4-ario obtuvo la menor mediana en 7 de 10 escalas sintéticas. La mediana de las razones 4-ario/binario fue 0.964 en el perfil disperso y 0.954 en el denso. Su menor profundidad teórica puede reducir el recorrido de \texttt{decrease\_key}; \texttt{extract\_min}, en cambio, compara hasta cuatro hijos por nivel. Los elementos están contiguos en el arreglo, pero no medimos fallos de caché. Además, los IQR de ambos montículos se solapan en las diez escalas: estos datos no bastan para atribuir la diferencia a la caché ni para afirmar una ventaja general del 4-ario.

\subsection{Desempeño del Montículo Binario en la Red Real de Cusco}
En la consulta OSM de Cusco (891 intersecciones y 2,139 aristas), el binario registró la menor mediana de cómputo: 15.65 ms, frente a 19.87 ms del 4-ario y 26.09 ms del Fibonacci. Se trata de un par origen--destino en una sola red; no prueba que el binario sea siempre más rápido por debajo de mil vértices. En particular, el 4-ario tuvo menor mediana en algunas escalas sintéticas menores que mil.

\subsection{Influencia de la Densidad del Grafo y Operaciones Decrease-Key}
En el recorte OSM se registran 2.40 aristas dirigidas por vértice. Para $|V|=5000$ disperso, el banco contabilizó 6,695 decrementos de clave entre 10,415 intentos de relajación (64.3\%). Este recuento depende de la instancia; por sí solo no explica la diferencia de tiempos entre colas.

\subsection{Evaluación de la Poda Espacial con Búsqueda Heurística \texorpdfstring{$A^*$}{A*}}
En el mismo par OSM, $A^*$ expandió 823 vértices frente a 887 de Dijkstra (7.2\% menos) y devolvió el mismo costo de viaje, 387.197 s. No se midió el tiempo de cómputo de $A^*$ como parte de la comparación de colas; reducir expansiones no demuestra por sí solo una reducción de tiempo. Tampoco se evaluaron pendientes ni pasajes peatonales en el \textit{snapshot} vial.

\subsection{Análisis de Sensibilidad de Factores Ambientales e Invariancia de Ruta}
La simulación con el factor uniforme $m_{\text{clima}} = 1.0691$ conserva el orden de costos entre caminos:
\begin{equation}
    w_{\text{clima}}(e) = m_{\text{clima}} \cdot w(e), \quad \forall e \in E \implies \sum_{e \in P_1} w_{\text{clima}}(e) < \sum_{e \in P_2} w_{\text{clima}}(e) \iff \sum_{e \in P_1} w(e) < \sum_{e \in P_2} w(e)
    \label{eq:invariancia_clima}
\end{equation}
El costo calculado pasó de 387.197 s a 413.952 s y la ruta de 17 vértices permaneció igual. Es un escenario uniforme basado en una observación meteorológica guardada, no una medición de tráfico ni una estimación calibrada por tramo.

\subsection{Guía Práctica y Marco de Decisión para Selección de Colas}
Para esta implementación en Python, las mediciones sugieren lo siguiente:
\begin{enumerate}
    \item \textbf{Binario:} Punto de partida sencillo; tuvo la menor mediana en la consulta OSM y en tres escalas sintéticas.
    \item \textbf{4-ario:} Conviene incluirlo al repetir mediciones en otras redes; registró menor mediana en siete escalas sintéticas, con IQR solapados frente al binario.
    \item \textbf{Fibonacci:} Conserva el interés por su cota amortizada, aunque no ganó ninguna de las escalas medidas aquí. No extrapolamos este resultado a C++ ni a otros tipos de grafo.
    \item \textbf{$A^*$:} Permite explorar menos vértices en la consulta OSM evaluada sin cambiar el costo; sus tiempos requieren una medición aparte.
\end{enumerate}

\subsection{Limitaciones del Estudio y Amenazas a la Validez}
\begin{itemize}
    \item \textbf{Intérprete y lenguaje:} Los tiempos se midieron en Python 3.12.8 sobre un AMD Ryzen 5 5600H. Otros equipos e intérpretes pueden dar resultados distintos.
    \item \textbf{Ausencia de mediciones físicas de memoria:} Las comparaciones espaciales se fundamentan en las cotas teóricas de punteros y estructuras de datos, al no haberse instrumentado un perfilador residente de asignación a nivel de sistema operativo.
    \item \textbf{Escenarios ambientales uniformes:} Los multiplicadores climáticos y de congestión operan como factores de sensibilidad homogéneos; futuros estudios requerirán calibrar coeficientes empíricos diferenciados por segmento vial a partir de aforos vehiculares en campo.
\end{itemize}
